% TOP_PROBLEM: {"id": 2307020, "title": "Research Problems in Function Theory — Problem 7.20", "category_id": 8, "category": {"id": 8, "name": "analysis", "display_name": "Analysis", "description": "Limits, continuity, calculus, and function theory.", "slug": "analysis", "order_index": 8, "created_at": "2026-07-31T15:26:25.671Z"}, "status": "open"}
% FIRST_POSTED: null
% ENTRY_KIND: statement_only
% Imported from: batch05.tex; UnsolvedMath ID 2307020
\documentclass[11pt]{article}
\usepackage{amsmath,amssymb,mathtools}
\usepackage{fontspec,unicode-math}
\setmainfont{Latin Modern Roman}
\setmathfont{Latin Modern Math}
\usepackage{xurl,hyperref}
\newcommand{\sref}[2]{\hyperlink{source:#1:#2}{[S#2]}}
\newcommand{\cataloguescope}[1]{\paragraph{Mathematical statement.}#1}
\begin{document}
% UnsolvedMath ID 2307020; source catalogue code AMR-022-7020
\setcounter{section}{2307019}
\section{Research Problems in Function Theory — Problem 7.20}\label{2307020}
{\small\sffamily UnsolvedMath ID 2307020 · Analysis}
\subsection{Definitions and mathematical statement}

\paragraph{Shared notation.}$\mathbb N=\{1,2,\ldots\}$, and $\mathbb Z,\mathbb Q,\mathbb R,\mathbb C$ denote the integers, rationals, reals and complex numbers. Any other conventions are specified locally.


Let $\mathbb D=\{z\in\mathbb C:|z|<1\}$, and let
\[X=\{(z_1,z_2):|z_1|\le1,\ |z_2|\le2\}\setminus\{(\zeta,\zeta):|\zeta|=1\}.
\]
A plurisubharmonic function is upper semicontinuous on its open domain and subharmonic on each complex line, with the value $-\infty$ permitted. Consider a function $F:X\to[-\infty,\infty)$ whose restriction to $\mathbb D^2$ is plurisubharmonic. No additional boundary continuity is imposed. At $z_1=z_2$ interpret $\log(1/|z_1-z_2|)$ as $+\infty$.

\paragraph{Statement recorded in the supplied source.}
Suppose that
\[F(z_1,z_2)\le\log\frac1{|z_1-z_2|}\quad((z_1,z_2)\in X),\]
and $F(z_1,z_2)\le0$ whenever $|z_1|=|z_2|=1$ and $z_1\ne z_2$. Must $F\le0$ throughout $\mathbb D^2$? \sref{2307020}{2}
The original domain includes $|z_2|\le2$; that printed domain is retained. Its boundary values and its interior plurisubharmonic restriction are stated separately. \sref{2307020}{2}
\subsection{Short English statement}
Does the given logarithmic bound and the nonpositive values on the distinguished boundary force a plurisubharmonic function to be nonpositive inside the bidisc?
\subsection{Sources}
\begin{description}
\item[\hypertarget{source:2307020:1}{S1}] ULAM.AI and the UnsolvedMath Contributors, \emph{UnsolvedMath: A Curated Collection of Open Mathematics Problems}, record 2307020 (AMR-022-7020). CC BY4.0. \url{https://huggingface.co/datasets/ulamai/UnsolvedMath}.
\item[\hypertarget{source:2307020:2}{S2}] W. K. Hayman and E. F. Lingham, \emph{Research Problems in Function Theory} (2018), Chapter7, Problem7.20 and its complete update. \url{https://arxiv.org/abs/1809.07200}.
\end{description}
\label{end:2307020}
\end{document}
